\scanpage{84}
\mspage{90}

Ainsi $f:X\longrightarrow S$ qui est une $\mathbb D$-fibration (propre et lisse) satisfait bien les conditions \inserted{\unsure{équiv.}} b), b') de la prop. \unsure{2,|2$_0$} (page 79), mais non la condition a), puisque le changement de base

\diagram{\begin{tikzcd}[ampersand replacement=\&]
X'=X\times_S\unsure{X|V|X'}\arrow[r,"q=\mathrm{pr}_1"]\arrow[d,"f'=\mathrm{pr}_2"'] \&X\arrow[d,"f"]\\
S'=X\arrow[r,"p=f"']\&S
\end{tikzcd}}

est une $\mathbb D$-équivalence $S'\to S$, alors que $X'\to X$ ne l'est pas.

Voici un autre contre'exemple, sur le même principe mais un peu plus simple. Soit $G$ un groupe \deleted{\unsure{fini}} \inserted{fini} (p. ex. $G=\mathbb Z/2\mathbb Z$), faisons le opérer sur une \deleted{catégorie} objet $\unsure{Z|\mathbb Z|X}$ de $\mathit{Cat}$, \unsure{soit|d'où} $S=[G]=B_G$ la catégorie réduite à un objet \inserted{$a$}, avec $G=\mathrm{End}(a)$.

\deleted{Alors on a}

\[
\deleted{f:X\longrightarrow S}
\]

\deleted{je dis que c'est une $\mathbb D$-fibration, propre et lisse, quel que soit $\mathbb D$. \unsure{Pour ceci|Bien plus}, considérons le chgt de base}

\[
\deleted{e\longrightarrow S}
\]

\deleted{\unsure{où}} Soit $E_G$ la catégorie \deleted{\unsure{$(B_G)_{/a}$}} \deleted{\unsure{connue de}} (\unsure{isom} : \deleted{\unsure{$(B_G)_{/a}$}}\inserted{$(B_G)_{/a}$})

\[
\mathrm{Ob}\,E_G=G,\qquad \mathrm{Hom}_{E_G}(g,g')=\{h\in G\mid\deleted{\unsure{h}}g=g'h\},
\]

composition évidente --- donc $E_G$ est un groupoïde trivial (\unsure{0|1}-connexe). Le groupe $G$ \deleted{\unsure{y}} opère à \deleted{droite}\inserted{gauche} librement, \inserted{par translations}, et $B_G=E_{\deleted{\unsure{G}}G}/G$.

% archive sheet number: 83
